% page 481 du fac-simile (image p-480 du scan)
% bloc 167.6 x 246.3 mm ; marge gauche 34.4 mm
\pgc{12.700mm}{0.000mm}{
\l{\cel{51}473}
\l{}
\l{\sou{pupilo}. (\sou{anat.)}\cel{2}Aperturo, meze di iriso, tra qua la lumo-radii}
\l{\cel{3}eniras la okulo. - DEFIRS.}
\l{}
\l{\sou{pupitro}. Moblo en qua existas surfaco plana, inklinita, sur\cel{1}\sou{qua}}
\l{\cel{3}onu pozas libri, muziko-kayeri, papero-folii, por lektar,}
\l{\cel{3}kantar, plear, skribar komfortoze. - FIS.}
\l{}
\l{\sou{pura.}\cel{2}I. Kun la naturo di qua ne esas mixita elemento stranjera.}
\l{\cel{3}- La\cel{1}\sou{matematiki pura, la mekaniko pura}\cel{1}: qui havas kom studi-}
\l{\cel{3}ajo nur la teorio, sen apliki praktikala. - II. Di qua la}
\l{\cel{3}naturo ne koruptesas da elemento stranjera. - DEFIS.}
\l{}
\l{\sou{pureo}.\cel{2}Disho ek legumi reduktita a paplo. - DEFIS.}
\l{}
\l{\sou{purgar}. (\sou{trans.}) I. Kuracar per eskapigar la feko quan kontenas}
\l{\cel{3}la intestini, per absorbigar lino-semini, agar-agaro, ricin-}
\l{\cel{3}oleo, aquo minerala sulfoza, e c. - II. (\sou{ulo)}. Liberigar de}
\l{\cel{3}to quo sordidigas, desparigas, koruptas, domajas (lu). - DEFIS.}
\l{}
\l{\sou{purgatorio}.\cel{2}(\sou{religio katol.})\cel{2}Loko ube la homi qui mortas ye}
\l{\cel{3}stando di graco, expiacas la peki pri qui li ne penitencis}
\l{\cel{3}sat multe en ica mondo. - DEFIS.}
\l{}
\l{\sou{puritano}. (\sou{kristanismo)}.Membro di sekto protestanta tre rigoroza,}
\l{\cel{3}qua aspiras ad igar kristanismo ri-esar pura quale lor lua}
\l{\cel{3}origino. - DEFIRS.}
\l{}
\l{\sou{purpuro.}\cel{2}Kolorizivo bele-reda quan antique onu obtenis de}
\l{\cel{3}L.\cel{1}\sou{murex trunculus}\cel{1}o\cel{1}\sou{murex brancari}s, molusko gasteropoda,}
\l{\cel{3}tipo di la tribuo \textquotedbl{}muricinei\textquotedbl{} qua vivas en la mari di qui}
\l{\cel{3}la aquo esas kolda o tepida, o fosila de la epoko \textquotedbl{}kretacea\textquotedbl{}. -}
\l{\cel{3}DEFIRS.}
\l{}
\l{\sou{puso}. (\sou{patol}.)Humoro morbala, densa, quan produktas la korupteso}
\l{\cel{3}di la tisui, lor lia inflameso. - EFIS.}
\l{}
\l{\sou{pustulo}. (\sou{patol.}) Mikra tumoro qua pusifas. - DEFIRS.}
\l{}
\l{\sou{putativa}. (yuro-cienco). Pri qua onu kredas ke lu esas valida}
\l{\cel{3}legale. - EFIS.}
\l{}
\l{\sou{puteo.}\cel{2}Exkavajo facita en sulo por tirar la aquo quan donas,}
\l{\cel{3}ye ta o ca profundeso-grado, fonti od infiltruri. -\cel{1}\sou{Puteo}}
\l{\cel{3}\sou{arteza}\cel{1}: Boro-kanaleto, facita en sulo per sondilo, qua}
\l{\cel{3}atingas (multa-kaze) profundeso-grado tre granda, til aquo-}
\l{\cel{3}strato. - FIS.}
\l{}
\l{putoro.\cel{2}(\sou{zool.}) Mamifero karnivora, ek la familio \textquotedbl{}martri\textquotedbl{},}
\l{\cel{3}qua odoras fetide. - L.\cel{1}\sou{mustela putorius}. - F.}
\l{}
\l{putrar.\cel{2}(netrans.) (\sou{pri, precipue, animalo o planto mortinta)}.}
\l{\cel{3}Deskompozar su. EFIS.}
\l{}
\l{}
\l{\cel{24}----}
}
